\section{Zusammenfassung und Ausblick}

Das Projekt \mira hat erfolgreich demonstriert, dass die Kombination aus Transfer Learning, gezielter Datensatz-Fusion und Post-Training-Quantisierung zu einem praxistauglichen, ressourceneffizienten KI-System für die automatische Wertstoffsortierung führt. Die vier zentralen Hypothesen wurden allesamt bestätigt:
\begin{itemize}
    \item \textbf{H1:} Transfer Learning mit MobileNetV2 übertraf das Scratch-CNN um +26,42 Prozentpunkte.
    \item \textbf{H2:} INT8-Quantisierung reduzierte die Modellgröße um den Faktor 9 auf \SI{2.61}{\mega\byte} bei vernachlässigbarem Genauigkeitsverlust (-0,07 pp).
    \item \textbf{H3:} Das YOLO11n-Detektionsmodell erreicht eine CPU-Inferenzrate von \SI{21.8}{FPS}, was für einen kontinuierlichen Sortierprozess ausreicht.
    \item \textbf{H4:} Der EMA-Filter ($\alpha=0,15$) eliminiert erfolgreich Signalflackern ohne kritischen Verzug.
\end{itemize}

Der entscheidende Durchbruch gelang durch die Erkenntnis, dass Data-Centric AI (Datenqualität und -vielfalt) die Modellleistung stärker beeinflusst als reine Architektur-Optimierungen. Der GIGO-Effekt beim Auto-Labeling zwang zur Adaption professioneller Datensätze, was die Generalisierungsfähigkeit dramatisch verbesserte.

Für eine industrielle Umsetzung sind mehrere Entwicklungsschritte notwendig:

\begin{enumerate}
    \item \textbf{Hardware-Integration:} Der Prototyp muss auf der Zielhardware (Raspberry Pi Zero 2W) implementiert und getestet werden. Dies beinhaltet die Anbindung an Servomotoren und Sensoren sowie die Optimierung der Inferenzpipeline für die ARM-Architektur.
    
    \item \textbf{Erweiterung der Klassifikation:} Das aktuelle fünfklassige System (Glas, Metall, Papier, Plastik, Restmüll) sollte um weitere relevante Abfallfraktionen erweitert werden, wie z.\,B. Biomüll, Elektroschrott oder Textilien. Dies erfordert die Integration neuer, spezialisierter Datensätze.
    
    \item \textbf{Synthetische Datengenerierung:} Um das Problem des Klassenungleichgewichts, insbesondere bei der Klasse \textit{Trash}, zu lösen, sollte ein Pipeline zur synthetischen Datengenerierung mittels Diffusionsmodellen (z.\,B. Stable Diffusion) etabliert werden. Dies würde es ermöglichen, realistische, annotierte Trainingsbilder für unterrepräsentierte Klassen zu erzeugen.
    
    \item \textbf{Online-Lernen:} Ein Mechanismus zum kontinuierlichen Lernen aus Fehlsortierungen im Live-Betrieb würde die langfristige Robustheit des Systems erhöhen. Hierfür müsste ein Feedback-Loop implementiert werden, der falsch klassifizierte Objekte in den Trainingsdatensatz zurückführt.
    
    \item \textbf{Energieeffizienz-Optimierung:} Für den batteriebetriebenen Einsatz sollte die Inferenzlatenz weiter reduziert werden, z.\,B. durch Pruning, Knowledge Distillation oder die Nutzung spezialisierter Edge-AI-Chips (z.\,B. Google Coral TPU).
\end{enumerate}

Die Kosten für ein solches System könnten bei einer Serienproduktion auf unter 100 Euro gesenkt werden, was es zu einer kosteneffizienten Alternative zu teuren industriellen Sortieranlagen macht. Insgesamt zeigt diese Arbeit, dass Deep Learning nicht nur in Cloud-Umgebungen, sondern auch auf extrem ressourcenbeschränkter Hardware praktikable Lösungen für reale Probleme liefern kann.